\chapter{Introduction}
\label{chap:introduction}

\section{Motivation}

Modern ray tracing GPUs -- NVIDIA RTX, Apple~GPU, Imagination~PowerVR --
deliver unprecedented rendering performance.
Yet every component of their ray tracing pipeline is proprietary:
the ISA, the BVH traversal unit, the shader execution model, the silicon.
This impedes research in three ways:
\begin{enumerate}
  \item \textbf{Non-reproducibility.} Algorithmic results depend on
    undocumented hardware behaviour.
  \item \textbf{Non-teachability.} Students learn to use a GPU, not
    to understand or build one.
  \item \textbf{Non-extensibility.} Novel architectural ideas cannot be
    prototyped without access to hardware design tools and foundries.
\end{enumerate}

This thesis addresses all three problems by building the entire stack
from first principles and releasing every component as open source.

\section{Contributions}

The main contributions of this thesis are:

\begin{enumerate}
  \item \textbf{\gonzales{}}: A production-quality, single-author
    path tracer capable of rendering the Disney Moana island scene.
    (\autoref{chap:gonzales})

  \item \textbf{\borg{}}: A Tile-Based Deferred Rendering GPU,
    implemented in Chisel, synthesized on FPGA, and taped out in
    GlobalFoundries 180\,nm silicon.
    (\autoref{chap:borg})

  \item \textbf{Ray tracing ISA extensions} for \borg{},
    enabling BVH traversal and intersection testing on hardware.
    (\autoref{chap:isa})

  \item \textbf{Gonzales-on-Borg}: The first demonstration of a
    production path tracer running on fully open-source hardware.
    (\autoref{chap:gonzales_on_borg})

  \item A comprehensive \textbf{evaluation} against CPU (Embree) and
    GPU (NVIDIA RTX) reference implementations.
    (\autoref{chap:evaluation})
\end{enumerate}

\section{Thesis Outline}

\autoref{chap:background} establishes the theoretical and architectural
background.
\autoref{chap:gonzales} presents the \gonzales{} renderer.
\autoref{chap:borg} describes the \borg{} GPU architecture.
\autoref{chap:isa} introduces the ray tracing ISA extensions.
\autoref{chap:gonzales_on_borg} presents the integrated system.
\autoref{chap:evaluation} evaluates performance and correctness.
\autoref{chap:conclusion} concludes and outlines future work.
